\scanpage{45}
\mspage{\unsure{1'|19}}
% archive foliation: 47
\uline{Corollaire 3} Soit $\underset{\text{ouvert}\atop\text{0-connexe}}{X}\subset X'$, $X'$ hol. 0-connexe. Alors $X$ est hyperbolique, sauf dans les cas suivants :

a) $X'$ est $\simeq\mathbb P^1_{\mathbb C}$ (resp. $\mathbb E^1_{\mathbb C}$, resp. $\mathbb G_{m\mathbb C}$), et $\operatorname{card}(X\setminus X')\leq2$ (resp. $\leq1$, resp. $=0$)

b) $X'$ est compacte de genre 1 et $X=X'$.

\uline{Remarque} Soit \deleted{\unsure{e}} $X$ une surface conforme 0-connexe (pas néc. orientable). On dit encore que $X$ est hyperbolique si $\widetilde X\simeq\mathbb H$, ou ce qui revient au $=$, si tte application conforme $\mathbb C\longrightarrow X$ est constante.

\margin{ou encore si le rev. \deleted{\unsure{absolu}} \inserted{des} orientations (qui est une surf. holomorphe) est hyperbolique.}

Ainsi les surfaces conformes \inserted{ss. bord,} 0-connexes \uline{non} hyperboliques sont, en plus des quatre cas énumérés dans (i) ci-dessus, les quotients de telles surfaces par des anti-involutions sans pts fixes. On trouve donc de plus
